\documentclass[11pt]{article}
\usepackage[margin=1in]{geometry}
\usepackage{amsmath,amssymb,amsthm}
\usepackage{hyperref}

\newtheorem{theorem}{Theorem}
\newtheorem{lemma}[theorem]{Lemma}
\newtheorem{proposition}[theorem]{Proposition}
\newtheorem{corollary}[theorem]{Corollary}
\theoremstyle{definition}
\newtheorem{definition}[theorem]{Definition}
\theoremstyle{remark}
\newtheorem{remark}[theorem]{Remark}

\title{Disproof of Conjecture 00000003481:\\
Thue numbers of trees are at most $4$, and $\pi(T)\le 2\log\Delta(T)$ fails}
\author{jilint777}
\date{2026-10-04}

\begin{document}
\sloppy
\maketitle

\begin{abstract}
Conjecture 00000003481 claims that the Thue number of the binary tree is four, that the Thue
number of general trees is unbounded, that $2\log\Delta$ is an upper bound for the Thue number of
a tree of maximum degree $\Delta$, and that this bound is attained asymptotically by subtrees of
the complete binary tree. We show that three of the four clauses are false. (i)~The bound
$\pi(T)\le 2\log\Delta(T)$ fails for every logarithm base $b>1$. The single edge $K_2$ has
$\Delta=1$ and $\pi=2>0=2\log 1$. Even for $\Delta\ge2$, the path $P_4$ has $\pi=3>2\log_b2$ for
every $b>4^{1/3}$, including bases $2$, $e$ and $10$. (ii)~The Thue number of trees is
\emph{bounded}: $\pi(T)\le 4$ for every tree (Bre\v{s}ar, Grytczuk, Klav\v{z}ar, Niwczyk and
Peterin, 2007). We include a complete proof of this bound that uses only Thue's theorem.
(iii)~Since $\pi\le4$ while $2\log\Delta\to\infty$, the bound is never attained asymptotically.
Moreover, subtrees of the complete binary tree have $\Delta\le3$, so they cannot form a family
with $\Delta\to\infty$. The refutation of (i) is formalised in Lean~4 (core library only).
So is the bound $\pi(T)\le4$ for every finite tree, and hence the refutation of (ii), with
Thue's theorem as the only explicit hypothesis. Graphs, paths, square factors, Thue colourings,
the Thue number and trees are all defined there from scratch. An independent Python script computes $\pi$ for all $987$ trees with at most $12$
vertices. It also finds that the complete binary tree of depth $6$ has $\pi=4$, which is
consistent with the remaining clause ``the Thue number of the binary tree is four''.
\end{abstract}

\section{The claim}

The conjecture (English version): \emph{``Definition: Thue coloring requires the color sequence
of any path to contain no square factor. Conjecture: The Thue number of the binary tree is four
and the Thue number of general trees is unbounded; the Thue-number upper bound for trees is twice
the logarithm of the maximum degree, attained asymptotically by subtree families of the complete
binary tree.''} The Chinese version has the same content.

It is a conjunction of four clauses:
\begin{itemize}
\item[(C1)] the Thue number of the binary tree is $4$;
\item[(C2)] the Thue number of general trees is unbounded;
\item[(C3)] $\pi(T)\le 2\log\Delta(T)$ for every tree $T$, where $\Delta(T)$ is the maximum degree;
\item[(C4)] the bound in (C3) is attained asymptotically by a family of subtrees of the complete
binary tree.
\end{itemize}
We show that (C2), (C3) and (C4) are false (Theorems~\ref{thm:c3}, \ref{thm:four},
Corollaries~\ref{cor:c2}, \ref{cor:c4}). One false clause already makes the conjunction false.
Section~\ref{sec:readings} lists the readings that each refutation covers. Clause~(C1) is not
refuted: Section~\ref{sec:c1} gives what we verified about it.

\section{Definitions}

All graphs are simple. A \emph{path} in a graph $G$ is a sequence $(v_1,\dots,v_\ell)$,
$\ell\ge1$, of pairwise distinct vertices in which $v_i v_{i+1}$ is an edge for every $i$. Any path
counts, between any two vertices; for trees these are not only root-to-leaf paths. A word $s$ has
a \emph{square factor} if $s=axxb$ for words $a,x,b$ with $x$ nonempty. A colouring $c$ of
$V(G)$ is a \emph{Thue colouring} (also called \emph{nonrepetitive}) if no path has a colour
sequence $c(v_1)\cdots c(v_\ell)$ with a square factor. The \emph{Thue number} $\pi(G)$ is the
least number of colours of a Thue colouring~\cite{AGHR}. A Thue colouring is in particular proper (take
$x$ of length~$1$), and $\pi$ is monotone under taking subgraphs. A word is \emph{square-free} if
it has no square factor.

\section{\texorpdfstring{Clause (C3) is false: $K_2$ and $P_4$}{Clause (C3) is false: K2 and P4}}

For a base $b>1$ and integers $k\ge0$, $\Delta\ge1$,
\begin{equation}\label{eq:log}
k\le 2\log_b\Delta \iff b^k\le\Delta^2 ,
\end{equation}
because $t\mapsto b^t$ is increasing. For a rational base $b=p/q$ ($p>q\ge1$) the right-hand side
reads $p^k\le \Delta^2q^k$. A larger base makes the bound smaller. So the claim for a base $b$
implies the claim for every base $b'\in(1,b]$. Every real $b>1$ lies above some rational
$p/q>1$, so refuting all rational bases $p/q>1$ refutes all real bases $b>1$.

\begin{lemma}\label{lem:small}
$\pi(K_2)=2$ and $\pi(P_4)=3$.
\end{lemma}

\begin{proof}
$K_2$: one colour gives the colour sequence $aa$ on the path through the edge. Two distinct
colours work, because the only paths are single vertices and the edge.

$P_4=v_1v_2v_3v_4$: suppose there is a Thue colouring with two colours. Then
$c(v_1)\ne c(v_2)\ne c(v_3)\ne c(v_4)$, because otherwise one edge carries the square $aa$. With
two colours this forces $c(v_1)=c(v_3)$ and $c(v_2)=c(v_4)$, so the whole path has colour sequence
$abab$, a square. Conversely, the colouring $1,2,3,1$ is a Thue colouring. Every path of $P_4$ is
a factor of $1231$ or of its reverse. Their factors of length $2$ have distinct letters, and the
only factor of length $4$ is $1231$ (or $1321$), which is not of the form $xx$.
\end{proof}

\begin{theorem}\label{thm:c3}
For every base $b>1$, the inequality $\pi(T)\le2\log_b\Delta(T)$ fails for some tree:
\begin{itemize}
\item $T=K_2$: $\Delta=1$ and $\pi=2>0=2\log_b1$, for every $b>1$;
\item $T=P_4$: $\Delta=2$ and $\pi=3>2\log_b2$ for every $b>4^{1/3}\approx1.587$; for example,
$2\log_22=2$, $2\ln2\approx1.386$ and $2\log_{10}2\approx0.602$;
\item $T=$ the complete binary tree $B_6$ of depth $6$ (computer search,
Section~\ref{sec:c1}): $\Delta=3$ and $\pi=4>2\log_b3$ for every $b>\sqrt3$; for example,
$2\log_23\approx3.17$.
\end{itemize}
\end{theorem}

\begin{proof}
Use Lemma~\ref{lem:small} and~\eqref{eq:log}. $b^2\le1$ is false for $b>1$, and $b^3\le4$ is
false for $b>4^{1/3}$.
\end{proof}

The last item shows that even the family named in (C4) violates the bound in (C3). The bound
also fails when it is rounded up: $\lceil 2\log_b1\rceil=0<2$, and $\lceil2\log_22\rceil=
\lceil2\ln2\rceil=2<3$. It fails as well for the Thue \emph{choice} number $\pi_{ch}\ge\pi$.

\section{\texorpdfstring{Every tree has $\pi(T)\le4$}{Every tree has pi(T) <= 4}}\label{sec:four}

This is a theorem of Bre\v{s}ar, Grytczuk, Klav\v{z}ar, Niwczyk and Peterin~\cite{BGKNP}; see
also Wood's survey~\cite[Thm.~4.1]{Wood}. For completeness we give a full proof. It uses the
following classical theorem.

\begin{theorem}[Thue 1906~\cite{Thue}]\label{thm:thue}
There is an infinite square-free word $w=w_0w_1w_2\cdots$ over $\{1,2,3\}$.
\end{theorem}

\begin{lemma}[K\"undgen--Pelsmajer~\cite{KP}; {\cite[Lemma~3.3]{Wood}}]\label{lem:kp}
Let $w$ be a square-free word over $\{1,2,3\}$. Let $u$ be the word obtained by inserting the
letter $4$ after every second letter,
$u=w_0w_1\,4\,w_2w_3\,4\,w_4w_5\,4\cdots$. Then (i) any three consecutive letters of $u$ are
pairwise distinct, and (ii) $u$ is square-free.
\end{lemma}

\begin{proof}
(i) Three consecutive letters of $u$ have one of the forms $w_{2j}w_{2j+1}4$,
$w_{2j+1}4w_{2j+2}$ or $4w_{2j+2}w_{2j+3}$. The two letters of $w$ in each are consecutive in $w$,
so they differ, since $w$ has no factor $aa$. Neither of them equals $4$.

(ii) Let $xx$ be a factor of $u$. Delete all letters $4$ from it. The result is a factor of $w$
of the form $x'x'$, where $x'$ is $x$ with its letters $4$ deleted. This is because the two copies
of $x$ have their $4$'s in the same positions. If $x'$ is empty, then $x$ consists of $4$'s only,
so $u$ has the factor $44$, which contradicts~(i). Otherwise $x'x'$ is a square in $w$, which is
a contradiction.
\end{proof}

\begin{lemma}[folding]\label{lem:fold}
Let $u$ be a square-free word in which any three consecutive letters are pairwise distinct. Let
$m,d_1,d_2\ge0$ with $m+\max(d_1,d_2)<|u|$ (or $u$ infinite). Then the word
\[
s=u_{m+d_1}u_{m+d_1-1}\cdots u_{m+1}\,u_m\,u_{m+1}\cdots u_{m+d_2}
\]
is square-free.
\end{lemma}

\begin{proof}
Index $s$ by $t\in\{-d_1,\dots,d_2\}$, so that $s_t=u_{m+|t|}$. Suppose $s$ has a square factor
occupying the positions $i,\dots,i+2k-1$ ($k\ge1$). Then $s_t=s_{t+k}$ for $i\le t\le i+k-1$.

If $i\ge0$, the factor lies in the part $t\ge0$, where $s_t=u_{m+t}$. So it is a square factor of
$u$, a contradiction. If $i+2k-1\le0$, it lies in the part $t\le0$, which is the reverse of a
factor of $u$. The reverse of a square is a square, so this is again a contradiction. So
$i<0<i+2k-1$.

Reversing $s$, i.e.\ replacing $t$ by $-t$ and swapping $d_1$ and $d_2$, maps squares to squares
and fixes the position $0$. If $0$ lies in the second half $[i+k,i+2k-1]$ of the square, then
after reversal it lies in the first half. So we may assume $i<0\le i+k-1$, i.e.\
$a:=-i\in[1,k-1]$.

\emph{Case $a\le k-2$.} The positions $-1,0,1$ lie in the first half $[i,i+k-1]$. Hence
$s_{k-1}=s_{-1}=u_{m+1}$ and $s_{k+1}=s_{1}=u_{m+1}$. Since $k-1\ge1>0$, we have
$s_{k-1}=u_{m+k-1}$ and $s_{k+1}=u_{m+k+1}$. So $u_{m+k-1}=u_{m+k+1}$, which contradicts the
hypothesis on three consecutive letters.

\emph{Case $a=k-1$} (so $k\ge2$). The square occupies $[1-k,k]$, and for $t\in[1-k,0]$ we have
$u_{m-t}=s_t=s_{t+k}=u_{m+t+k}$. Put $j=-t\in[0,k-1]$ and $y_j=u_{m+j}$ for $0\le j\le k$. Then
$y_j=y_{k-j}$, so $y_0\cdots y_k$ is a palindrome of length $k+1\ge3$. If $k$ is even, then
$y_{k/2-1}=y_{k/2+1}$. If $k$ is odd, then $y_{(k-1)/2}=y_{(k+1)/2}$. Both contradict the
hypothesis.
\end{proof}

\begin{theorem}[\cite{BGKNP}]\label{thm:four}
Every tree $T$, finite or infinite, satisfies $\pi(T)\le4$.
\end{theorem}

\begin{proof}
Let $u$ be the infinite word of Lemma~\ref{lem:kp} built from Thue's word
(Theorem~\ref{thm:thue}). Fix a root $r$, and colour every vertex $v$ with $u_{\mathrm{dist}(r,v)}$.
Every non-root vertex has exactly one neighbour of smaller depth, its parent, and every edge joins
a vertex to its parent. Let $P=(v_0,\dots,v_\ell)$ be a path. An interior vertex $v_j$ cannot have
both $v_{j-1}$ and $v_{j+1}$ as its parent, because they are distinct. So the depth sequence along
$P$ has no interior local maximum. It changes by exactly $1$ at each step, so it first decreases
by steps of $1$ down to the depth $m$ of the top vertex, and then increases by steps of $1$. Hence
the colour sequence of $P$ is exactly the word $s$ of Lemma~\ref{lem:fold}, with $d_1,d_2$ the
lengths of the two monotone parts. That word is square-free.
\end{proof}

For a finite tree only a finite prefix of $u$ is needed. For finite trees, this proof is
formalised in Lean, with Thue's theorem as an explicit hypothesis (Section~\ref{sec:lean}). Its existence for every length is still
Thue's theorem. Bre\v{s}ar et al.~\cite{BGKNP} also show that $\pi(T)=4$ for some trees, so the
bound is tight. Section~\ref{sec:c1} gives an explicit example.

\begin{corollary}\label{cor:c2}
Clause (C2) is false: $\sup_T \pi(T)=4<\infty$ over all trees.
\end{corollary}

\begin{corollary}\label{cor:c4}
Clause (C4) is false. Let $(T_n)$ be any family of trees with $\Delta(T_n)\to\infty$. Then
$\pi(T_n)/(2\log\Delta(T_n))\le 4/(2\log\Delta(T_n))\to0$, so the bound $2\log\Delta$ is not
attained asymptotically, in any base. Moreover, a subtree of the complete binary tree has maximum
degree at most $3$. So subtree families of the complete binary tree have bounded $\Delta$, and an
asymptotic statement in $\Delta$ about them is impossible. Every such subtree satisfies
$\pi\le 4$, while $2\log_b\Delta\le 2\log_b 3$ is fixed.
\end{corollary}

\section{\texorpdfstring{What we know about clause (C1)}{What we know about clause (C1)}}
\label{sec:c1}

Let $B_h$ be the complete binary tree of depth $h$. Its root has degree $2$, and its other
internal vertices have degree $3$. An exhaustive backtracking search (\texttt{verify.py}) gives
\[
\pi(B_1)=2,\quad \pi(B_2)=\pi(B_3)=\pi(B_4)=\pi(B_5)=3,\quad \pi(B_6)=4.
\]
The colourings found for $h\le5$ are re-checked on all paths. The non-existence of a Thue
$3$-colouring of $B_6$ (127 vertices) is confirmed by two different searches: breadth-first order
with colour-symmetry breaking, and depth-first order without it. By monotonicity and
Theorem~\ref{thm:four}, $\pi(B_h)=4$ for all $h\ge6$, and also for the infinite binary tree. So
(C1) holds for the infinite or sufficiently deep complete binary tree, and we do not refute it.
Under the reading ``every binary tree has Thue number $4$'' it would be false, because $P_4$ is a
binary tree with $\pi=3$. We do not rely on this reading. The result $\pi(B_6)=4$ is a computer
result and is not formalised in Lean.

\section{Readings covered}\label{sec:readings}

The Lean result \texttt{conjecture\_00000003481\_false} refutes the literal (C3). Under the
asymptotic reading of (C3), the disproof goes through (C2), whose refutation, $\pi(T)\le4$ for
every tree, is proved completely in Section~\ref{sec:four}. That refutation is also formalised in Lean,
modulo Thue's theorem as an explicit hypothesis (Section~\ref{sec:lean}).

\begin{itemize}
\item \textbf{(C3) as a universal bound $\pi(T)\le2\log_b\Delta(T)$} (the literal reading).
False for every base $b>1$ by $K_2$. False for every $b>4^{1/3}$, in particular $b=2,e,10$, by
$P_4$ even when $\Delta\ge2$ is required, and false for $b>\sqrt3$ by $B_6$. Also false with
$\lceil\cdot\rceil$, and for the choice number $\pi_{ch}\ge\pi$.
\item \textbf{(C3) only ``for large $\Delta$'' or up to $O(1)$.} Then (C3) is true but trivial,
by Theorem~\ref{thm:four}: $4\le2\log_b\Delta$ as soon as $\Delta\ge b^2$. In this case (C4)
is false (Corollary~\ref{cor:c4}), and so is (C2). The conjunction is false under every reading.
\item \textbf{(C2)} is false under every reading of ``Thue number'' as the chromatic-type
parameter defined in the statement (Corollary~\ref{cor:c2}), for finite and for infinite trees.
\emph{Disclosure:} for the different parameter $\pi_{ch}$ (nonrepetitive \emph{list}
colouring), trees do have unbounded values (Fiorenzi, Ochem, Ossona de Mendez and
Zhu~\cite{FOOZ}). The statement defines colourings, not list colourings, so this does not apply.
Even for $\pi_{ch}$, the literal (C3) fails by $K_2$ and $P_4$.
\end{itemize}

\section{Formal verification (Lean 4)}\label{sec:lean}

The directory \texttt{lean4/} contains a Lean~4.19.0 project that uses only the core library.
\texttt{Main.lean} defines, from scratch:
\begin{itemize}
\item \texttt{Graph} (vertices $0,\dots,n-1$, Boolean adjacency);
\item \texttt{IsPath} (nonempty, duplicate-free, consecutive vertices adjacent; all paths, not
only root-to-leaf ones);
\item \texttt{HasSquare s} ($s=a\,x\,x\,b$ with $x\ne[\,]$);
\item \texttt{Nonrepetitive}, \texttt{Colorable G k} and \texttt{IsThueNumber G k} (the least
such $k$, proved unique);
\item \texttt{IsTree}: a nonempty simple graph in which any two vertices are joined by exactly
one path (the textbook definition; for $K_2$ and $P_4$ we also check ``connected with $n-1$
edges''), and \texttt{maxDeg};
\item \texttt{LogBound p q $\Delta$ k} $:=p^k\le\Delta^2q^k$, which is~\eqref{eq:log} for the
base $p/q$.
\end{itemize}
The main results are:
\begin{itemize}
\item \texttt{thue\_K2}: $\pi(K_2)=2$, and \texttt{thue\_P4}: $\pi(P_4)=3$;
\item \texttt{K2\_tree} and \texttt{P4\_tree}: both are trees, with \texttt{maxDeg} $1$ and $2$;
\item \texttt{conjecture\_00000003481\_false}: for all $p>q$, $\neg$\,\texttt{Claim p q}, where
\texttt{Claim p q} says that every tree with $\Delta\ge1$ and Thue number $k$ satisfies
\texttt{LogBound p q $\Delta$ k};
\item \texttt{claimCol\_false}: the same for the colouring form ``some Thue colouring with
$\le2\log\Delta$ colours exists'', which does not use the minimum;
\item \texttt{claimDeg2\_false}: the claim restricted to $\Delta\ge2$ is false whenever
$p^3>4q^3$, for example for bases $2$ and $27/10<e$;
\item \texttt{claim\_mono}: a larger base gives a stronger claim;
\item \texttt{nonvacuous}: the hypotheses are satisfiable, and Thue colouring is a nontrivial
condition.
\end{itemize}

\paragraph{Clause (C2) in Lean (Section 7 of \texttt{Main.lean}).} The hypothesis is
\[
\texttt{ThueHyp} := \forall L,\ \exists w:\mathbb N\to\mathbb N,\ (\forall i,\ w_i<3)\ \wedge\
\forall i\,k,\ 0<k\to i+2k\le L\to \exists j<k,\ w_{i+j}\ne w_{i+j+k},
\]
i.e.\ Thue's theorem~\cite{Thue,Berstel} in finite form: a square-free ternary word of every
length. The list form \texttt{ThueHypList} (lists of every length with entries $<3$ and
$\neg$\,\texttt{HasSquare}) implies it (\texttt{thueHyp\_of\_list}).
\texttt{thueHyp\_40} checks the hypothesis for $L=40$ by \texttt{decide}. From it, Lean
proves:
\begin{itemize}
\item \texttt{kp\_sqfree}, \texttt{kp\_ne1}, \texttt{kp\_ne2}: Lemma~\ref{lem:kp} for
$u_i=3$ if $i\equiv2\pmod 3$ and $u_i=w_{2\lfloor i/3\rfloor+(i\bmod 3)}$ otherwise;
\item \texttt{fold\_sqfree}: Lemma~\ref{lem:fold}, for $t\mapsto u_{m+|t-a|}$, with all six
positions of the turning point relative to the square;
\item \texttt{valley}: along any path of a graph with a layering with unique parents, the depth
is $m+|t-a|$;
\item \texttt{tree\_layering}: every tree (in the sense of \texttt{IsTree}) has such a layering,
namely the depth $d(v)=|P_v|-1$, where $P_v$ is the unique path from $0$ to $v$;
\item \texttt{tree\_colorable\_four}: assuming \texttt{ThueHyp}, every tree has a Thue
$4$-colouring; hence $\pi(T)\le4$ (\texttt{thueNumber\_le\_four});
\item \texttt{conjecture\_C2\_false}: assuming \texttt{ThueHyp}, it is not the case that for
every $k$ some tree $T$ satisfies $\neg$\,\texttt{Colorable T k}. The variant with $\pi$ is
\texttt{conjecture\_C2\_false'}; the variant assuming the list form is
\texttt{conjecture\_C2\_false\_list}.
\end{itemize}
These use the same \texttt{Graph}, \texttt{IsPath}, \texttt{HasSquare}, \texttt{Colorable} and
\texttt{IsTree} as the main theorem.
Nonrepetitiveness of the explicit colourings is decided over all candidate vertex sequences. The
completeness of this enumeration is proved: by a pigeonhole lemma, a path has at most $n$
vertices. Only the axioms \texttt{propext}, \texttt{Classical.choice} and \texttt{Quot.sound}
are used. There is no \texttt{sorry} and no \texttt{native\_decide}.

The following are \emph{not} in Lean, only here and in \texttt{verify.py}:
\begin{itemize}
\item Thue's theorem itself (Theorem~\ref{thm:thue}), which is an explicit hypothesis;
\item the case of infinite trees in Theorem~\ref{thm:four} (Lean graphs are finite);
\item Corollary~\ref{cor:c4};
\item the computation $\pi(B_6)=4$;
\item the passage from rational to real bases, which is the one-line monotonicity argument after
\eqref{eq:log}; its rational form is \texttt{claim\_mono}.
\end{itemize}

\section{Independent check (Python)}

\texttt{verify.py} uses only the standard library. It does the following:
\begin{itemize}
\item computes $\pi(K_2)=2$ and $\pi(P_4)=3$ by enumerating all colourings and all paths;
\item generates all trees with $n\le12$ vertices up to isomorphism ($987$ trees; the counts match
OEIS A000055) and computes $\pi$ for each by backtracking, re-checking each colouring on all
paths;
\item finds the maximum to be $3$ for $n\le12$, and counts the trees that violate (C3): $10$ for
base $2$ and $661$ for base $e$;
\item computes $\pi(B_h)$ for $h\le6$, with the two independent searches for $B_6$. Both
searches are fast: each shows that $B_6$ (127 vertices) has no Thue $3$-colouring in well under
a second (about $0.03$\,s). The whole script runs in about $7$\,s;
\item checks the construction of Theorem~\ref{thm:four}. It builds a ternary square-free word,
applies the K\"undgen--Pelsmajer insertion and checks the depth colouring on all $11006$ rooted
trees with $n\le12$ and on $B_1,\dots,B_8$.
\end{itemize}

\begin{thebibliography}{9}
\bibitem{Thue} A.~Thue, \"Uber unendliche Zeichenreihen, \emph{Norske Vid.\ Selsk.\ Skr.\ I
Mat.-Nat.\ Kl.\ Christiania} 7 (1906), 1--22.
\bibitem{AGHR} N.~Alon, J.~Grytczuk, M.~Ha{\l}uszczak, O.~Riordan, Nonrepetitive colorings of
graphs, \emph{Random Structures Algorithms} 21 (2002), 336--346.
\bibitem{BGKNP} B.~Bre\v{s}ar, J.~Grytczuk, S.~Klav\v{z}ar, S.~Niwczyk, I.~Peterin,
Nonrepetitive colorings of trees, \emph{Discrete Math.} 307 (2007), 163--172.
\bibitem{KP} A.~K\"undgen, M.~J.~Pelsmajer, Nonrepetitive colorings of graphs of bounded
tree-width, \emph{Discrete Math.} 308 (2008), 4473--4478.
\bibitem{FOOZ} F.~Fiorenzi, P.~Ochem, P.~Ossona de Mendez, X.~Zhu, Thue choosability of trees,
\emph{Discrete Appl.\ Math.} 159 (2011), 2045--2049.
\bibitem{Berstel} J.~Berstel, \emph{Axel Thue's papers on repetitions in words: a
translation}, Publications du LaCIM 20, Universit\'e du Qu\'ebec \`a Montr\'eal, 1995.
\bibitem{Wood} D.~R.~Wood, Nonrepetitive graph colouring, \emph{Electron.\ J.\ Combin.}, Dynamic
Survey DS24 (2021); arXiv:2009.02001.
\end{thebibliography}

\end{document}
